\section{Introduction}
\label{sec:introduction}

Scientific papers are the foundational medium for scientific knowledge and are now collected in massive digital collections~\citep{lo2020s2orc,priem2022openalex}. 
Embedding models transform each paper into a vector, enabling recommendation, large-scale scientific analysis~\citep{reimers2019sbert,cohan2020specter,su2023instructor,beltagy2019scibert,ostendorff2022scincl,fortunato2018science}, and predicting new discoveries~\citep{tshitoyan2019mat2vec,peng2021periodical2vec,krenn2020predicting,sourati2023accelerating}. 
However, most embeddings are not directly interpretable: embedding is an information compression process that retains information needed for downstream tasks and discards the rest. 
Because the current paradigm mainly focuses on similarity search, the embedding models are trained to preserve the similarity structure but obscures the actual idea represented by individual papers~\citep{weller2025limitations,kuratov2025cramming}.  
Since the original content is not recoverable, interpretating the idea of a given point---which may not be paired with an actual paper such as a center of a cluster---typically requires heuristics such as labeling the nearest documents~\citep{grootendorst2023keyllm,pham2024topicgpt,wang2023goalex} or inverting the vector to text~\citep{morris2023vec2text,song2020leakage,li2023geia}, both of which only approximate the underlying idea. Alternative approaches build interpretability into the vectors themselves by defining each coordinate as a sparse concept~\citep{koh2020concept,bhalla2024splice}, a fixed answer~\citep{benara2024qaemb}, or a feature from a sparse autoencoder~\citep{oneill2024sae}, but these methods limit interpretation to a predefined vocabulary.

\afterpage{%
\begin{figure}[t]
\centering
\includegraphics[width=\linewidth]{doc2lora-figs.pdf}
\caption{
\textbf{Doc2LoRA pipeline.}
\textbf{Encode:} The frozen \doctolora hypernetwork $f_{\mathrm{enc}}$ compresses a document into an embedding tensor $\mathcal{V}$ of shape $(n_L, n_R, n_d)$, sitting between the MLP and LoRA heads (\secref{doc2lora-embedding}). A bijective map $g_\theta$ (\secref{kron-method}) aligns this space for downstream search.
\textbf{Decode:} Any point, including interpolations and averages, is mapped back via $g_\theta^{-1}$ and expanded by $f_{\mathrm{dec}}$ into a LoRA update $\Delta \matrix{W}$, enabling natural language queries to the base model. 
}
\label{fig:method}
\end{figure}%
}

We propose a new embedding paradigm where any point in an embedding space represents a large language model (LLM) to which we can query about the idea represented by that point.
This conversational embedding enables idea arithmetic such as averaging a cluster of ideas to form a category-level description or interpolating between two ideas for novel ideas, an operation motivated by the long-standing finding that new knowledge arises from recombining existing findings~\citep{swanson1986undiscovered,uzzi2013atypical,foster2015tradition,youn2015invention,shibayama2021novelty}.
Because every point in this space is an LLM, the result of arithmetic operations is also an LLM we can query about the resulting idea. 

We enable this paradigm by leveraging Doc-to-LoRA~\citep{sakana2024doc2lora}, which extracts a low-dimensional latent representation for each document. While large language models are typically parameterized by billions of weights, \doctolora generates a compact adapter that captures only the document-specific changes while shared prior knowledge remains in the frozen base model and is not encoded in the embedding.
We use the compact latent tensor produced by its encoder-decoder architecture as our document embedding. 
This means each embedding isolates what makes that document unique, omitting redundant commonalities present across all documents. 
Although \doctolora was originally designed for context distillation~\citep{snell2022context}, we show that its latent space has an intuitive geometric structure and is effective as an embedding.

While several works represent texts such as task instruction and documents as LoRA adapters~\citep{chen2025ga,su2025parametricrag,ye2021hypter,mahabadi2021hyperformer,phang2023hypertuning,charakorn2025texttolora}, we select \doctolora~\citep{sakana2024doc2lora} as our base hypernetwork for its crucial properties: identifiability (the same document always maps to the same adapter) and scalability (efficient embedding of large corpora). 
For example, Parametric RAG~\citep{su2025parametricrag}, which fine-tunes a LoRA adapter per document, lacks identifiability due to random initialization, resulting in different adapters for the same input~\citep{garipov2018loss,frankle2020linear,entezari2022permutation,ainsworth2023gitrebasin}. GenerativeAdapter~\citep{chen2025ga} achieves identifiability and fast, fixed-size adapter generation, but produces prohibitively large adapters, e.g., $268$M parameters ($537$\,MB in \texttt{bf16}) per document for Mistral-7B—making corpus-level indexing impractical. 
In contrast, \doctolora creates a compact, consistent adapter per document in one pass, compressing each to a $147$k-dimensional latent tensor (and poolable to $18$k for search; see \secref{doc2lora-embedding}), achieving both identifiability and scalability.

We tested the effectiveness of \doctolora embeddings in terms of idea arithmetic along with similarity search.
On idea arithmetic, we show that \doctolora embeddings support operations such as averaging a cluster of ideas into a category-level description or interpolating between two ideas to mix them (\secref{cluster-labeling}, \secref{fusion}).
We further compare against \texttt{T2L}~\citep{charakorn2025texttolora}, a hypernetwork that also emits LoRA adapters from a frozen \texttt{GTE} encoder~\citep{li2023gte} rather than from a coordinate system learned as a document representation (\secref{results}).
On similarity search, tested on next-paper prediction, collaboration, topic classification, and author-name disambiguation, a small bijective transform trained on OpenAlex citation pairs lifts the generation-trained geometry to a level comparable with \texttt{SPECTER2}~\citep{singh2023scirepeval}, \texttt{Instructor}~\citep{su2023instructor}, \texttt{EmbeddingGemma}~\citep{embeddinggemma2025}, and \texttt{GTE}~\citep{li2023gte} (\secref{similarity}).
The transform is linear and invertible, so the embedding remains decodable.
\texttt{SBERT}~\citep{reimers2019sbert} remains the strongest encoder overall, though \doctolora with $g_\theta$ performs competitively on all tasks.

We present the \doctolora latent as an embedding space where each point corresponds to a unique, decodable LLM, allowing idea retrieval, arithmetic, and synthesis. Using the public SakanaAI \doctolora checkpoint as a fixed hypernetwork, we introduce two architectural advances: (1) treating the post-MLP tensor as a linearly equivariant, mean-pooled index for compact, composable representation (\secref{doc2lora-embedding}), and (2) a bijective $262$k-parameter retrieval map that maintains adapter decodability (\secref{kron-method}). This framework enables practical idea arithmetics and semantic operations, improving interpretability and flexibility over prior embedding methods (\secref{results}).

%これまでの研究では，タスクごとに共有ベースモデルをファインチューニングした際に生じる重みをタスクの表現として用いてきた．これらの重みは，精度向上のために平均化したり~\citep{wortsman2022soups}，スキルの追加・削除を担うタスクベクトルとして合成されたり~\citep{ilharco2023task}，学習年次をまたいで補間されたりしている~\citep{nylund2024time}．1タスクでファインチューニングしたモデル同士は，低損失領域で連結されることが示されており~\citep{gueta2023knowledge}，個別に学習されたアダプタ同士は再利用可能なモジュールとして統合できる~\citep{chronopoulou2023adaptersoup,huang2024lorahub,yadav2023ties}．また，重み自体をデータとして読み出すこともでき，ネットワークの精度を重みから予測したり~\citep{unterthiner2020predicting}，モデル動物園内で重みそのものから表現を学習することも可能である~\citep{schurholt2021hyper}．
%
%しかし，ファインチューニングによるアダプタは，コーパススケールで文書埋め込みとして利用することはできない．同じ文書を個別にファインチューニングした場合，それぞれの重みは異なる順序関連のバシンに落ち込み，互いに比較できなくなる~\citep{}．また，文書ごとにファインチューニングを施す手法は，1回読むよりも1桁ほど多くの計算量がかかり，文書1つにつき数MBの重みを保存する必要があるため，スケールしない~\citep{su2025parametricrag}．

%これに対し，ハイパーネットワークを用いることで，各文書を1回のフォワードパスで，1つの共有座標系上のユニークなアダプタに写像できる~\citep{ye2021hypter,mahabadi2021hyperformer,phang2023hypertuning,chen2025ga,charakorn2025texttolora}．

%We represent a paper by the low-rank adapter (LoRA)~\citep{hu2022lora} that a hypernetwork~\citep{ha2017hypernetworks} generates from it (\figref{method}).
%The language model with this adapter can answer questions about the paper without seeing the paper in the prompt.
%Because every point of the space expands into such an adapter, averages, interpolations, and other constructed vectors are likewise interpretable.
%We decode a paper by expanding its embedding into a LoRA and querying the frozen base model (\secref{decode}).
%The same vectors remain usable for similarity search.
%After a small bijective transform trained on OpenAlex citation pairs, cached \doctolora embeddings overtake \texttt{SPECTER2}~\citep{singh2023scirepeval} on next-paper prediction and \texttt{SPECTER} on four of the five author-name disambiguation datasets~\citep{subramanian2021s2and} (\secref{similarity}).
